% ৩.১৫ সার্বজনীন গেট
\section{Universal Gates: NAND, NOR}

\subsection{NAND}
\begin{itemize}
\item $Y=\overline{A\cdot B}$; NAND দিয়ে AND, OR, NOT তৈরি
\end{itemize}

\subsection{NOR}
\begin{itemize}
\item $Y=\overline{A+B}$; NOR দিয়ে সব গেট নির্মাণ
\end{itemize}

\begin{learningbox}
এই টপিক শেষে শিক্ষার্থী পারবে:
\begin{itemize}
\item NAND ও NOR gate-এর truth table লিখতে।
\item universal gate কেন বলা হয় ব্যাখ্যা করতে।
\item NAND/NOR দিয়ে NOT, AND, OR gate নির্মাণের ধারণা দেখাতে।
\end{itemize}
\end{learningbox}

\begin{definitionbox}{Universal gate}
যে gate একাই ব্যবহার করে মৌলিক gateসহ যেকোনো logic circuit তৈরি করা যায়, তাকে universal gate বলা হয়। NAND ও NOR universal gate।
\end{definitionbox}

\begin{center}
\begin{tabular}{|c|c|c|c|}
\hline
$A$ & $B$ & NAND $\overline{A\cdot B}$ & NOR $\overline{A+B}$ \\
\hline
0 & 0 & 1 & 1 \\
\hline
0 & 1 & 1 & 0 \\
\hline
1 & 0 & 1 & 0 \\
\hline
1 & 1 & 0 & 0 \\
\hline
\end{tabular}
\end{center}

\begin{center}
\fbox{\parbox{0.9\textwidth}{
\centering
\textbf{Circuit placeholders}\\[4pt]
NAND দিয়ে NOT: দুই input একসাথে $A$।\\
NAND দিয়ে AND: NAND output-কে আবার NAND inverter-এ দাও।\\
NOR দিয়ে NOT: দুই input একসাথে $A$।\\
NOR দিয়ে OR: NOR output-কে আবার NOR inverter-এ দাও।
}}
\end{center}

\begin{boardbox}{বোর্ড ফোকাস}
Universal gate প্রমাণের প্রশ্নে অন্তত NOT, AND এবং OR gate তৈরি দেখাতে হয়। শুধু ``সব gate তৈরি করা যায়'' লিখলে পূর্ণ নম্বর নাও আসতে পারে।
\end{boardbox}
